\documentclass[12pt,a4paper]{report}

% ============================================================
%  CODIFICACIÓN E IDIOMA
% ============================================================
\usepackage[utf8]{inputenc}
\usepackage[T1]{fontenc}
\usepackage[spanish]{babel}

% ============================================================
%  MATEMÁTICAS
% ============================================================
\usepackage{amsmath}
\usepackage{amssymb}

% ============================================================
%  GRÁFICOS Y TABLAS
% ============================================================
\usepackage{graphicx}
\usepackage{booktabs}
\graphicspath{{imagenes/}}

% ============================================================
%  BIBLIOGRAFÍA (natbib + bibtex)
% ============================================================
\usepackage[authoryear,round]{natbib}
\bibliographystyle{plainnat}

% ============================================================
%  HIPERVÍNCULOS (cargar al final del preámbulo)
%  plainpages=false + pdfpagelabels=true evitan el conflicto de destinos
%  duplicados que genera hyperref al pasar de numeración romana a arábiga.
% ============================================================
\usepackage[hidelinks,plainpages=false,pdfpagelabels=true]{hyperref}

% ============================================================
%  DATOS DE PORTADA
%  Título e institución tomados del artículo base del proyecto
%  ("Mitigación de Sesgos en IA: Análisis del Sistema de Admisión
%  Escolar en Chile"). Esta tesis tiene un único autor.
%  \gradoTesis y \directorTesis quedan como placeholders porque no
%  se cuenta con esa información confirmada — complétala tú.
% ============================================================
\newcommand{\tituloTesis}{Mitigación de Sesgos en IA: Análisis del Sistema de Admisión Escolar en Chile}
% \autorTesis se usa en la portada (texto plano).
\newcommand{\autorTesis}{Diego Villegas Cisternas}
\newcommand{\institucionTesis}{Universidad Técnica Federico Santa María}
\newcommand{\facultadTesis}{Departamento de Informática}
\newcommand{\carreraTesis}{[Nombre de la carrera]}       % PLACEHOLDER — completar
\newcommand{\gradoTesis}{[Grado académico o título]}     % PLACEHOLDER — completar
\newcommand{\directorTesis}{[Nombre del profesor(a) guía]} % PLACEHOLDER — completar
\newcommand{\fechaTesis}{\today}
\newcommand{\ciudadTesis}{Santiago, Chile}

\title{\tituloTesis}
\author{Diego Villegas Cisternas}
\date{\fechaTesis}

\begin{document}

% ============================================================
%  PORTADA + ÍNDICE (numeración romana)
%  Todo el preliminar (portada + índice) se envuelve con pageanchor=false
%  para que hyperref no genere destinos "page.N" que luego choquen con los
%  de la numeración arábiga cuando el contador de página se reinicia
%  (evita el warning de "destino duplicado").
% ============================================================
\begingroup
\hypersetup{pageanchor=false}

\begin{titlepage}
    \centering
    \vspace*{1cm}

    {\large \institucionTesis \par}
    {\large \facultadTesis \par}

    \vspace{2.5cm}

    {\Huge \bfseries \tituloTesis \par}

    \vspace{2.5cm}

    {\large \autorTesis \par}

    \vspace{1.5cm}

    {\normalsize Tesis presentada para optar al \gradoTesis \ en \carreraTesis \par}

    \vspace{1.5cm}

    {\normalsize Profesor(a) guía: \directorTesis \par}

    \vfill

    {\large \ciudadTesis \par}
    {\large \fechaTesis \par}
\end{titlepage}

\pagenumbering{roman}
\tableofcontents
\clearpage

\endgroup
\pagenumbering{arabic}

% ============================================================
%  CAPÍTULOS
%  (cada archivo contiene únicamente un \chapter{...} placeholder)
% ============================================================
\chapter{Resumen}

En el contexto del rápido aumento del uso de la inteligencia artificial en el diseño de experiencias de usuario, surgen inquietudes importantes sobre la falta de transparencia y la presencia de sesgos en los sistemas algorítmicos que toman decisiones de alto impacto social. Esta memoria aborda esos problemas a partir de un caso real: el Sistema de Admisión Escolar (SAE) de Chile, que asigna estudiantes a establecimientos educacionales mediante un algoritmo de Aceptación Diferida que, si bien es formalmente correcto y está regido por la Ley de Inclusión Escolar N.\textsuperscript{o}~20.845, no comunica su lógica a las familias que lo usan, generando una percepción extendida de arbitrariedad.

El trabajo se apoya en un diagnóstico heurístico de calidad web del SAE real —realizado por el equipo del proyecto Fondecyt N.\textsuperscript{o}~1250492 de la Universidad de Chile, que reveló un cumplimiento de 51\,\% en el sitio informativo y 61\,\% en la plataforma de postulación— y en una revisión sistemática de 96 publicaciones académicas sobre transparencia algorítmica con evaluación de usuarios finales (2017--2025). A partir de ese diagnóstico y de esa revisión se definió un perfil de usuario objetivo (Daniela González) y se construyó, en sucesivas iteraciones documentadas entre mayo y julio de 2026, un prototipo de interfaz que aplica principios de divulgación progresiva, explicabilidad contextualizada y controles interactivos, evolucionando desde un prototipo HTML autocontenido hasta una aplicación web funcional construida con React.

Una verificación interna de trazabilidad frente a la matriz de problemas del diagnóstico registra un cumplimiento del 100\,\% de los puntos aplicables a nivel de código. Una revisión adicional del sitio oficial del SAE, realizada como parte de este trabajo, mostró que solo dos de las seis brechas prioritarias identificadas —búsqueda y encontrabilidad, e inclusión parcial— han sido cerradas de forma independiente por el sistema real, mientras que la transparencia algorítmica interactiva y la interoperabilidad con ClaveÚnica permanecen sin resolver, lo que confirma la vigencia del problema abordado.

\textbf{Esta memoria se encuentra en una etapa intermedia de desarrollo.} El prototipo aún no ha sido sometido a una reevaluación heurística formal ni a una prueba de usabilidad con usuarios reales; ambas validaciones están diseñadas como un plan de cinco fases (Capítulo~3) pero su ejecución, junto con la discusión y las conclusiones que de ella se deriven, queda como trabajo pendiente. Este documento presenta, por tanto, el diagnóstico, el marco teórico, la metodología completa y el estado actual del desarrollo, dejando explícitamente abiertos los capítulos de Resultados de validación, Discusión y Conclusiones hasta que dicha evidencia esté disponible.

\chapter{Introducción}

El avance acelerado y la creciente integración de la inteligencia artificial (IA) en diversos sectores han transformado el diseño y la experiencia de productos y servicios digitales, desde la salud y las finanzas hasta la logística y la atención ciudadana \citep{strobel2019aspects}. Esta integración generalizada eleva el diseño ético de la experiencia de usuario (UX) de una característica deseable a una necesidad fundamental: la ética en UX implica tomar decisiones de diseño conscientes para beneficiar a los usuarios y actuar con responsabilidad social \citep{keepcoding2024etica}, evitando activamente los llamados ``patrones oscuros'' —elementos de interfaz que inducen a los usuarios a tomar decisiones que no los benefician \citep{checkealos2024etica}.

Un componente central de ese diseño ético es la \emph{transparencia algorítmica}: la capacidad de un sistema de IA para que su proceso de toma de decisiones, la lógica subyacente y los resultados generados sean interpretables y comprensibles para las personas \citep{ibm2024transparencia,transparenciaalgoritmica}. La ausencia de transparencia no es un problema meramente estético; puede derivar en decisiones incorrectas, filtraciones de información, discriminación por raza o género, y una pérdida generalizada de confianza institucional \citep{unir2024sesgo,skillnest_sesgo,agoro2019case}.

En Chile, el \textbf{Sistema de Admisión Escolar (SAE)} es un caso paradigmático de esta tensión. El SAE es la plataforma del Ministerio de Educación que asigna estudiantes a establecimientos públicos y particulares subvencionados mediante un algoritmo basado en el mecanismo de \emph{Aceptación Diferida} (\emph{Deferred Acceptance}), formulado originalmente por \citet{galeshapley1962} y adaptado al problema de asignación ``muchos a uno'' que caracteriza la admisión escolar. En Chile, este mecanismo está regido por la Ley de Inclusión Escolar N.\textsuperscript{o}~20.845 y considera criterios de prioridad definidos legalmente: hermanos matriculados en el establecimiento, pertenencia al 15\,\% de estudiantes prioritarios, hijos de funcionarios del colegio y condición de exalumno \citep{mineduc_sae,epstein2017mecanismos}.

El algoritmo del SAE es, en ese sentido, formalmente correcto: implementa un mecanismo matemáticamente probado y aplicado en sistemas de admisión escolar de ciudades como Boston, Nueva York, Singapur y París, entre otros. Sin embargo, el sistema no comunica esta lógica a las familias que lo usan. El resultado es una percepción extendida de arbitrariedad —el SAE es descrito coloquialmente como una ``tómbola''— que erosiona la confianza pública incluso cuando los resultados de asignación son objetivamente justos según los criterios legales vigentes. Esta brecha fue confirmada empíricamente por la evaluación heurística de calidad web realizada por el equipo del proyecto Fondecyt N.\textsuperscript{o}~1250492 de la Universidad de Chile en marzo de 2026 \citep{moralesvargas2026informe}, que encontró un cumplimiento de apenas 51\,\% en el sitio informativo del SAE y 61\,\% en su plataforma de postulación, con la dimensión de \emph{transparencia y apertura} entre las peor evaluadas del sitio informativo (0\,\% de cumplimiento): no existe ninguna sección que explique públicamente el funcionamiento del algoritmo, sus criterios de asignación o sus resultados agregados.

Esta brecha es abordable desde el diseño: no requiere modificar el algoritmo ni la legislación vigente, sino intervenir en la capa de presentación e interacción para hacer visible una lógica que hoy permanece oculta para el usuario final. Es precisamente en ese punto donde converge la evidencia de la literatura internacional sobre transparencia algorítmica con evaluación de usuarios finales: enfoques como la divulgación progresiva, la explicabilidad contextualizada y los controles interactivos han demostrado, en decenas de estudios empíricos, que es posible aumentar la comprensión y la confianza del usuario sin exponer los detalles técnicos del algoritmo ni sobrecargar cognitivamente la interfaz.

\section*{Problema de investigación}
\addcontentsline{toc}{section}{Problema de investigación}

¿Es posible diseñar una interfaz de usuario que, sin alterar el algoritmo de asignación del SAE, mejore significativamente la comprensión, la confianza y la percepción de justicia de las familias respecto del proceso de admisión escolar, aplicando principios de transparencia algorítmica validados por la literatura científica?

\section*{Objetivo general}
\addcontentsline{toc}{section}{Objetivo general}

Diseñar, construir y evaluar un prototipo de interfaz de usuario que mejore la transparencia y la experiencia de usuario del Sistema de Admisión Escolar chileno, fundamentado en evidencia de la literatura sobre transparencia algorítmica y en un diagnóstico heurístico del sistema real.

\section*{Objetivos específicos}
\addcontentsline{toc}{section}{Objetivos específicos}

\begin{enumerate}
    \item Sistematizar la evidencia científica disponible sobre transparencia algorítmica con evaluación de usuarios finales, identificando qué enfoques de diseño funcionan y cuáles no.
    \item Diagnosticar, mediante un instrumento de evaluación heurística de calidad web, las brechas específicas del sitio informativo y de la plataforma de postulación del SAE.
    \item Diseñar un caso de estudio representativo y un perfil de usuario objetivo (persona) que guíen las decisiones de diseño del prototipo.
    \item Construir un prototipo de interfaz —primero como maqueta de alta fidelidad y posteriormente como aplicación web funcional— que aborde las brechas prioritarias detectadas, con énfasis en el módulo explicativo del algoritmo de asignación.
    \item Evaluar el prototipo mediante pruebas de usabilidad con usuarios reales, midiendo claridad, confianza y utilidad percibida.
    \item Discutir los resultados obtenidos a la luz de la literatura revisada y proponer recomendaciones para una eventual implementación institucional.
\end{enumerate}

\section*{Estructura de la tesis}
\addcontentsline{toc}{section}{Estructura de la tesis}

El Capítulo~2 presenta el marco teórico: los fundamentos de diseño ético en IA, la distinción entre transparencia, explicabilidad e interpretabilidad, los sesgos algorítmicos y los marcos regulatorios vigentes, además de la síntesis de la revisión sistemática de 96 publicaciones sobre transparencia algorítmica con evaluación de usuarios finales. El Capítulo~3 describe la metodología aplicada, en dos frentes complementarios: el diagnóstico heurístico del SAE real y el proceso de diseño, construcción y validación del prototipo. El Capítulo~4 reporta los resultados obtenidos, tanto del diagnóstico como de las pruebas de usabilidad. El Capítulo~5 discute estos resultados en relación con la literatura y con el estado actual del sitio oficial del SAE. Finalmente, el Capítulo~6 presenta las conclusiones y las líneas de trabajo futuro.

\chapter{Marco Teórico}

\section{Diseño ético en experiencias impulsadas por IA}

El diseño ético de UX se entiende como la práctica de tomar decisiones de diseño de manera consciente, con el objetivo de beneficiar a los usuarios y actuar con responsabilidad social \citep{keepcoding2024etica}. Un aspecto central de esta práctica es la prevención activa de los \emph{patrones oscuros} (\emph{dark patterns}): elementos de diseño implementados deliberadamente para inducir a los usuarios a realizar acciones no deseadas o tomar decisiones que no los benefician \citep{checkealos2024etica}. Estas prácticas no solo erosionan la confianza del usuario, sino que han motivado respuestas regulatorias concretas, como las multas de la \emph{Federal Trade Commission} de Estados Unidos a empresas como Epic Games y Amazon por el uso indebido de patrones oscuros \citep{checkealos2024etica}. El diseño ético, por contraste, prioriza el bienestar del usuario y garantiza que sus necesidades se atiendan de forma justa y respetuosa, en lugar de explotar vulnerabilidades cognitivas para obtener beneficios de corto plazo.

Esta idea conecta directamente con un principio más amplio de la disciplina de interacción humano-computador: el \textbf{diseño centrado en el usuario} (\emph{user-centered design}), que postula que las decisiones de un sistema interactivo deben fundamentarse en las necesidades, capacidades y contexto real de las personas que lo usarán, y no únicamente en la lógica interna del sistema o en la conveniencia de quien lo construye. Una de las herramientas más difundidas de este enfoque es la construcción de \emph{personas}: arquetipos de usuario, construidos a partir de datos reales de caracterización, que sintetizan un perfil representativo del segmento objetivo y sirven como referencia constante durante el proceso de diseño. En el Capítulo~3 se describe cómo esta técnica se aplicó concretamente en este trabajo mediante el arquetipo ``Daniela González''.

\section{Transparencia, explicabilidad e interpretabilidad}

La literatura distingue tres conceptos que suelen usarse de manera intercambiable pero que tienen alcances distintos \citep{ibm2024transparencia}:

\begin{itemize}
    \item \textbf{Transparencia}: se refiere al contexto y la metodología usada para desarrollar el modelo (qué datos se usaron, bajo qué criterios se diseñó).
    \item \textbf{Explicabilidad}: es la capacidad del sistema para justificar sus decisiones mediante explicaciones claras y contextualizadas a un caso específico.
    \item \textbf{Interpretabilidad}: es la capacidad del modelo para presentar una lógica de decisión comprensible en términos generales, más allá de un caso puntual.
\end{itemize}

El incumplimiento de estos tres aspectos intensifica desigualdades sociales existentes, al reforzar y amplificar sesgos previos: puede derivar en decisiones incorrectas, exposición de información sensible, y discriminación por raza o género \citep{unir2024sesgo}. Casos ampliamente documentados ilustran esta consecuencia: el sistema COMPAS en Estados Unidos, criticado por predicciones judiciales sesgadas, o el algoritmo BOSCO en España, que enfrentó desafíos legales al denegar beneficios sociales de forma no transparente \citep{legalprod,wikipedia_transparencia}.

\subsection{Sesgos algorítmicos}

Los sesgos algorítmicos pueden introducir desigualdades o reforzar estereotipos existentes, afectando negativamente a determinados grupos demográficos. La literatura distingue al menos tres mecanismos de origen \citep{skillnest_sesgo,ibm_iaresponsable}:

\begin{itemize}
    \item \textbf{Sesgo en los datos de entrenamiento}: surge cuando los datos usados para entrenar un modelo contienen prejuicios raciales, de género o socioeconómicos preexistentes.
    \item \textbf{Sesgo algorítmico propiamente tal}: se origina en el diseño o implementación del algoritmo, cuando quienes lo desarrollan magnifican, consciente o inconscientemente, sesgos ya existentes.
    \item \textbf{Sesgo por retroalimentación}: ocurre cuando resultados ya sesgados se reintroducen al sistema como nuevos datos de entrada, reforzando un ciclo de discriminación.
\end{itemize}

En el ámbito de la salud, por ejemplo, se ha documentado cómo sistemas de IA han reproducido sistemáticamente desigualdades de acceso y tratamiento médico, lo que subraya una responsabilidad ética compartida entre quienes diseñan y quienes regulan estos sistemas \citep{rivas2024transparency}. Estos casos son relevantes para el SAE porque ilustran que un algoritmo puede ser técnicamente correcto y, aun así, producir efectos percibidos como injustos si el proceso que lo rodea no es comprensible ni auditable por las personas afectadas.

\subsection{Marcos regulatorios}

En respuesta a estos desafíos, distintas entidades internacionales han desarrollado marcos normativos que enfatizan la transparencia y la equidad en sistemas de IA. Entre los más relevantes destacan el \emph{NIST AI Risk Management Framework 1.0}, que define prácticas de gestión de riesgo mediante cuatro funciones —gobernar, mapear, medir y gestionar— \citep{nist2023airmf}; la norma \emph{ISO/IEC 42001:2023}, que propone un sistema integral de gestión de IA centrado en gobernanza, gestión de riesgos y trazabilidad de las decisiones algorítmicas \citep{iso2023aims}; y los principios actualizados de la OCDE, que recalcan la transparencia, la robustez técnica, la equidad y la rendición de cuentas como pilares de una adopción sostenible de IA. Ninguno de estos marcos exige o prohíbe un algoritmo de asignación específico; todos, en cambio, ponen el énfasis en la \emph{comunicación} de su funcionamiento hacia las personas afectadas, que es exactamente el problema que aborda esta memoria.

\section{Revisión sistemática: transparencia algorítmica con evaluación de usuarios finales}

Como base empírica adicional para el diseño del prototipo, se sistematizó una revisión de \textbf{96 publicaciones académicas} sobre transparencia algorítmica que incluyen evaluación con usuarios finales, cubriendo el período 2017--2025 y consultando cuatro bases de datos (SciSpace, SciSpace Texto Completo, Google Scholar y ArXiv). Esta revisión —independiente del artículo base del proyecto— permite contrastar los hallazgos generales del campo con el caso específico del SAE.

\subsection{Qué funciona}

La evidencia converge de manera consistente en cinco enfoques de diseño efectivos:

\begin{itemize}
    \item \textbf{Divulgación progresiva}: mostrar primero información esencial y revelar el detalle técnico solo si el usuario lo solicita, evitando la fijación temprana en errores o detalles irrelevantes \citep{springerwhittaker2019,springerwhittaker2020}.
    \item \textbf{Explicabilidad contextualizada}: explicaciones ligadas a la situación específica del usuario son más efectivas que las descripciones abstractas del funcionamiento general del sistema, especialmente en contextos de decisión pública como trámites de visado o beneficios sociales \citep{nefedov2022}.
    \item \textbf{Controles interactivos con retroalimentación}: permitir que el usuario explore el sistema (simuladores, ajustes, comparadores) aumenta la sensación de agencia y reduce la percepción de arbitrariedad, siempre que existan etiquetas claras y retroalimentación inmediata \citep{kim2021recommender,feddersen2024forecasting}.
    \item \textbf{Presentación multidimensional}: combinar gráficos con cifras mejora la comprensión y la calidad de la decisión del usuario, frente a la presentación de un único indicador agregado \citep{glazerman2018multidim}.
    \item \textbf{Reportes y explicaciones co-diseñadas con usuarios no expertos}: los formatos de explicación diseñados junto a los propios usuarios finales —en lugar de derivados solo de consideraciones técnicas— se ajustan mejor a sus necesidades reales de información \citep{peng2023codesign,bobek2024xai}.
\end{itemize}

\subsection{Qué no funciona}

La literatura también documenta con igual claridad los enfoques que fallan o que pueden ser contraproducentes:

\begin{itemize}
    \item La \textbf{transparencia técnica sin filtro} —mostrar código, ecuaciones o el detalle interno completo del algoritmo— reduce la confianza en usuarios no técnicos en lugar de aumentarla \citep{springerwhittaker2019}.
    \item El \textbf{exceso de información en una sola pantalla} genera parálisis decisional y peores evaluaciones de satisfacción.
    \item Las \textbf{explicaciones globales} (``así funciona el sistema para todos'') son sistemáticamente menos efectivas que las \textbf{explicaciones locales} (``así funcionó en tu caso específico'').
    \item Otorgar a usuarios sin entrenamiento \textbf{control de grano fino} sobre componentes complejos del sistema puede producir ajustes peores que los del sistema automático \citep{feddersen2024forecasting}.
    \item Distintos métodos técnicos de explicación (por ejemplo, mapas de saliencia o atribuciones locales) pueden producir niveles de comprensión humana similares entre sí, lo que sugiere que la evaluación debería centrarse en el ajuste a la tarea del usuario y no solo en la novedad del método \citep{dawoud2023comparative}.
\end{itemize}

\subsection{Relevancia para el caso SAE}

El Sistema de Admisión Escolar presenta características que amplifican estos hallazgos: sus usuarios son, en su mayoría, apoderados con baja o media alfabetización digital; la decisión tiene un alto impacto emocional (la educación de un hijo o hija); y el algoritmo subyacente tiene aspectos contraintuitivos —por ejemplo, que la posición de un colegio en la lista de preferencias no garantiza por sí sola la asignación—. Esto convierte el diseño orientado a la transparencia en un requisito funcional, no en un elemento cosmético, y confirma la pertinencia de aplicar al SAE los mismos principios (divulgación progresiva, explicabilidad contextualizada, controles interactivos) que la literatura identifica como efectivos en otros dominios de alto impacto, como la administración pública \citep{nefedov2022}, la salud \citep{rivas2024transparency} y los sistemas de recomendación \citep{masciari2024systematic}.

\section{El algoritmo del Sistema de Admisión Escolar}

El SAE utiliza un algoritmo basado en el mecanismo de \emph{Aceptación Diferida} (\emph{Deferred Acceptance}, DA), propuesto originalmente por \citet{galeshapley1962} para el problema del emparejamiento estable entre dos conjuntos de agentes. Aunque su formulación original resolvía el problema del emparejamiento uno a uno, el modelo fue extendido a problemas de asignación \emph{muchos a uno}, que es precisamente el caso de la admisión escolar, donde cada colegio tiene capacidad para recibir a múltiples estudiantes.

El mecanismo opera de forma iterativa \citep{epstein2017mecanismos}:

\begin{enumerate}
    \item \textbf{Postulación}: las familias envían un listado ordenado de los colegios de su preferencia.
    \item \textbf{Primera ronda}: cada estudiante postula a su primera preferencia. Los colegios revisan a sus postulantes y, según sus criterios de prioridad, aceptan tentativamente a los estudiantes hasta llenar sus vacantes; los no aceptados quedan rechazados.
    \item \textbf{Rondas sucesivas}: los estudiantes rechazados en la ronda anterior postulan al siguiente colegio de su lista. Los colegios reevalúan a sus nuevos postulantes junto con los ya aceptados tentativamente, conservando siempre a los de mayor prioridad.
    \item \textbf{Finalización}: el proceso termina cuando todos los estudiantes han sido asignados a un establecimiento (o agotan su lista de preferencias).
\end{enumerate}

En Chile, este mecanismo está regido por la Ley de Inclusión Escolar N.\textsuperscript{o}~20.845 y considera cuatro criterios de prioridad: tener un hermano matriculado en el establecimiento, pertenecer al 15\,\% de estudiantes prioritarios, ser hijo o hija de un funcionario del colegio, y ser exalumno que no haya sido expulsado \citep{mineduc_sae}. Es importante notar que, si bien el SAE no tiene su código fuente públicamente disponible, existe documentación oficial y trabajos académicos —como la tesis de \citet{epstein2017mecanismos}— que permiten reconstruir su lógica interna y sus criterios de priorización. La falta de acceso completo al código limita, sin embargo, la posibilidad de una auditoría exhaustiva por parte de terceros, lo que refuerza la necesidad de una capa de comunicación explicativa dirigida al usuario final, que es precisamente el problema que aborda esta memoria.

\section{Accesibilidad web como condición de transparencia}

Un sistema no puede considerarse transparente si una parte relevante de sus usuarios no puede, en la práctica, acceder a su contenido. Por ello, este trabajo adopta como referencia las \emph{Web Content Accessibility Guidelines} (WCAG) del World Wide Web Consortium, en particular los criterios de nivel AA: contraste mínimo de 4.5:1 entre texto y fondo, posibilidad de navegar todo el sitio exclusivamente con teclado, texto alternativo descriptivo en imágenes funcionales, y tamaños de fuente legibles en dispositivos móviles. Estos criterios no son un añadido secundario en el contexto del SAE: el diagnóstico heurístico de la Universidad de Chile encontró una dimensión de \emph{inclusión} con 0\,\% de cumplimiento en el sitio informativo real, lo que confirma que la accesibilidad es, en este caso concreto, una de las brechas más urgentes y no solo un ideal de buenas prácticas.

\section{Prototipado iterativo asistido por herramientas de IA generativa}

Un aspecto metodológico propio de este trabajo —y que se detalla en el Capítulo~3— es que el prototipo se construyó mediante un proceso iterativo de \emph{research through design} apoyado en un agente de inteligencia artificial generativa (Claude, de Anthropic, operado a través de su interfaz de desarrollo Claude Code) como herramienta de programación asistida. Este tipo de flujo de trabajo, cada vez más habitual en el desarrollo de software y en la construcción de prototipos de investigación, no reemplaza las decisiones de diseño ni la evaluación crítica de quien dirige el proyecto: el agente ejecuta instrucciones de diseño explícitas, redactadas por el autor en documentos de especificación (descritos en el Capítulo~3), y el resultado de cada iteración fue revisado, corregido y redirigido antes de avanzar a la siguiente. Documentar este proceso con honestidad —incluyendo sus herramientas— es, en sí mismo, coherente con el principio de transparencia que esta memoria defiende para el propio SAE.

\chapter{Metodología}

Esta investigación emplea una metodología de \textbf{diseño aplicado} (\emph{research through design}), en la que el propio acto de diseñar, construir y refinar un artefacto —en este caso, un prototipo de interfaz— constituye una forma de generar conocimiento sobre el problema abordado. El proceso se organiza en cuatro fases: diagnóstico del sistema real, síntesis de evidencia investigativa, construcción iterativa del prototipo, y un plan de validación que, como se explicita en este capítulo, se encuentra actualmente en etapa de diseño y ejecución pendiente.

\section{Fase 1 — Diagnóstico heurístico del SAE}

La primera fase consistió en diagnosticar el estado real del SAE mediante el instrumento de evaluación heurística de calidad web desarrollado por SISIB (Universidad de Chile), aplicado en marzo de 2026 por el equipo del proyecto Fondecyt N.\textsuperscript{o}~1250492 \citep{moralesvargas2026informe} sobre dos componentes del sistema: el sitio web informativo (\texttt{sistemadeadmisionescolar.cl}) y la plataforma transaccional de postulación.

El instrumento combina veinte dimensiones generales de calidad web —contenido y lenguaje claro, usabilidad, accesibilidad, arquitectura de información, búsqueda y encontrabilidad, responsividad móvil, diseño e imagen institucional, seguridad, tecnología, atención a la ciudadanía, audiovisualidad, enfoque de género, inclusión, promoción, transparencia y apertura, prevención de errores, facilidad de acceso, interacción y retroalimentación, y rapidez de respuesta, entre otras— con un bloque adicional de catorce factores específicos para plataformas de selección escolar. Cada dimensión se descompone en indicadores de chequeo asociados a tres niveles de cumplimiento —\emph{imprescindible}, \emph{esperable} y \emph{deseable}—, y el puntaje final se calcula como el porcentaje de indicadores cumplidos sobre el total de indicadores aplicables, ponderado por nivel.

Este diagnóstico entregó el \emph{baseline} cuantitativo (51\,\% en el sitio informativo, 61\,\% en la plataforma de postulación) que se usó como insumo directo para priorizar las decisiones de diseño de la Fase~3, y que —según se detalla en la Sección~3.5— debe usarse también como referencia de comparación en la validación pendiente del prototipo.

\section{Fase 2 — Revisión sistemática de literatura}

En paralelo, se realizó una revisión sistemática de 96 publicaciones académicas sobre transparencia algorítmica con evaluación de usuarios finales (2017--2025), mediante búsqueda en cuatro fuentes —SciSpace (búsqueda semántica), SciSpace de texto completo, Google Scholar (consulta booleana con palabras clave dirigidas) y ArXiv (\emph{preprints} de ciencias de la computación e interacción humano-computador)—. Los resultados de las cuatro fuentes fueron fusionados, deduplicados y reordenados por relevancia, priorizando los estudios más pertinentes para el problema de investigación. Los criterios de inclusión favorecieron estudios que implementaran algún mecanismo de transparencia algorítmica \emph{y} lo evaluaran empíricamente con usuarios reales, descartando propuestas puramente teóricas sin evaluación. Esta revisión, descrita en el Capítulo~2, aportó los principios de diseño que fundamentan las decisiones de la Fase~3.

Como antecedente directo y motivación inicial de este trabajo, cabe mencionar que un estudio previo del mismo equipo autor —centrado específicamente en el módulo de explicación del resultado de asignación y construido como maqueta estática en Figma— fue sometido a una prueba de usabilidad con diez personas, con resultados favorables en comprensión y confianza percibida. Ese estudio, sin embargo, corresponde a un artefacto y a una instancia de evaluación \textbf{distintos} de los abordados en esta memoria: se cita aquí únicamente como antecedente que motivó la decisión de profundizar y ampliar el prototipo al resto del flujo del SAE, no como resultado propio de este trabajo. La evaluación del prototipo desarrollado en esta memoria (Fase~3) queda descrita como pendiente en la Sección~3.5.

\section{Fase 3 — Diseño y construcción iterativa del prototipo}

\subsection{Caso de estudio y persona}

Para guiar las primeras decisiones de diseño se definió un caso de estudio hipotético (un perfil de estudiante ficticio, una lista ordenada de preferencias de colegios, y un estado de vacantes definido para cada uno) y un arquetipo de usuario objetivo, \textbf{Daniela González}: 35 años, educación media completa, alfabetización digital básica-intermedia, acceso a internet principalmente por teléfono móvil, que postula a un máximo de tres colegios y valora la seguridad, la cercanía y el prestigio del establecimiento. Su frustración principal —documentada en el informe Fondecyt— es percibir el SAE como una ``tómbola'' que no explica sus reglas. De este perfil se derivaron implicancias de diseño concretas: interfaz mobile-first con 375~px como referencia primaria, nivel de lectura objetivo de 6.\textsuperscript{o} básico, flujo de postulación en un máximo de tres pasos, y explicaciones del algoritmo concretas y personalizadas al caso del usuario.

\subsection{Cronología real de las iteraciones}

A diferencia de un desarrollo de software convencional con un único entregable final, este prototipo se construyó como una sucesión de iteraciones documentadas, cada una guiada por un documento de especificación explícito redactado por el autor y ejecutada con apoyo de un agente de IA generativa como herramienta de programación asistida (ver Capítulo~2, Sección~2.6). La Tabla~\ref{tab:iteraciones} resume esta cronología, reconstruida a partir del historial de control de versiones del proyecto.

\begin{table}[h]
\centering
\caption{Cronología de iteraciones del prototipo (mayo--julio de 2026).}
\label{tab:iteraciones}
\small
\begin{tabular}{p{2.2cm}p{9.5cm}}
\toprule
\textbf{Fecha} & \textbf{Hito} \\
\midrule
Mayo (semana 1) & Redacción de \texttt{CLAUDE.md} (especificación base del proyecto) y construcción de la \textbf{Iteración~1}: prototipo HTML autocontenido de una sola página (\texttt{prototipo\_SAE\_mejora.html}), que cubre las cinco secciones mínimas exigidas por el diagnóstico (inicio, módulo de algoritmo, ficha de colegio, flujo de postulación, panel de seguimiento). \\
Mayo (semana 2) & Redacción de \texttt{CLAUDE\_v2.md} (especificación de mejoras) e incorporación de los documentos de investigación de soporte (revisión de literatura, informe heurístico, plan de metodología). Construcción de la \textbf{Iteración~2}: segunda versión HTML (\texttt{prototipo\_SAE\_v2.html}) que corrige limitaciones técnicas de la V1 (responsividad limitada a dos \emph{breakpoints}, buscador no funcional, simulador con lógica simplificada) y agrega funcionalidades nuevas: comparador de colegios, calendario interactivo, y una lógica de simulación más cercana al mecanismo de Aceptación Diferida. \\
Mayo (semana 3) & \textbf{Migración de arquitectura}: el prototipo HTML monolítico se reescribe como aplicación \textbf{React + Vite}, con el objetivo de pasar de un archivo estático a una arquitectura de componentes mantenible y extensible. El contenido textual de la versión HTML se extrae programáticamente para no perder el trabajo de redacción ya validado. \\
Mayo (semana 4) & Se incorpora \emph{shadcn/ui} como sistema de componentes base y se reconstruyen sobre React las páginas de algoritmo, inicio, postulación y seguimiento. \\
Junio (semana 3) & Se documentan formalmente las brechas del diagnóstico en \texttt{feedback\_sae\_problemas.md} y se construye \texttt{plan\_mejora\_sae.md}: una matriz de trazabilidad de aproximadamente 62 puntos que vincula cada problema detectado con la decisión de diseño correspondiente y su estado de implementación. Los prototipos HTML anteriores se archivan como \emph{baseline} histórico. \\
Junio (semana 4) & Se redacta \texttt{CONTEXTO\_CLAUDE\_CODE.md} como documento de continuidad del proyecto, y se implementa el \textbf{Comparador de colegios} (\texttt{ComparadorPage.jsx}), identificado como la funcionalidad significativa pendiente de mayor prioridad. \\
Julio (commit del 12/07) & Se incorporan un \textbf{tour guiado} de \emph{onboarding} (\texttt{GuidedTour.jsx}) y una reescritura sustancial del flujo de postulación (validación de RUT, análisis de probabilidad de asignación por colegio en la propia lista de preferencias). \\
Julio (trabajo posterior, aún sin registrar en control de versiones) & Se construye \texttt{NotasPage.jsx}: una página dedicada de ``Notas de diseño'' que documenta los seis principios de investigación aplicados en el prototipo, cada uno con su referencia bibliográfica, el hallazgo empírico que lo respalda y los componentes concretos donde se implementó, cerrando así el requisito de trazabilidad visible exigido por \texttt{CLAUDE.md}. Se enlaza la nueva ruta \texttt{/notas} desde el pie de página del sitio. \\
\bottomrule
\end{tabular}
\end{table}

\emph{Nota:} la última fila de la Tabla~\ref{tab:iteraciones} corresponde a cambios ya presentes en el código del prototipo pero que, a la fecha de esta memoria, aún no han sido registrados como \emph{commit} en el historial de control de versiones del proyecto.

Este historial es relevante metodológicamente por dos razones. Primero, confirma que el prototipo no es un artefacto estático sino el resultado de sucesivas correcciones guiadas por el propio diagnóstico y por la literatura revisada, lo que es coherente con el carácter iterativo del \emph{research through design}. Segundo, dado que cada iteración fue ejecutada con apoyo de un agente de IA a partir de instrucciones explícitas del autor, este trabajo documenta también —de forma honesta y verificable a través del propio historial de control de versiones— el uso de esa herramienta como parte del proceso metodológico, y no únicamente el resultado final.

\subsection{Estado actual del prototipo}

Al cierre de la última iteración registrada, el prototipo cubre las cinco secciones exigidas por el diagnóstico original (inicio, módulo del algoritmo con simulador funcional, ficha de colegio, flujo de postulación en tres pasos con inicio de sesión simulado vía ClaveÚnica, y panel de seguimiento), además de tres funcionalidades adicionales no exigidas explícitamente por el diagnóstico pero incorporadas durante el desarrollo: un comparador de colegios, un tour guiado de \emph{onboarding}, y una página de ``Notas de diseño'' que documenta explícitamente la trazabilidad entre cada principio de la revisión de literatura (Capítulo~2) y su implementación concreta en el código. Una verificación interna de trazabilidad frente a la matriz de \texttt{plan\_mejora\_sae.md} registra 61 de 62 puntos abordados en el código —cifra que, con la incorporación de la página de notas de diseño, alcanza en la práctica el 100\,\% de los puntos de esa matriz—; se trata, sin embargo, de una \textbf{autoevaluación de implementación}, no de una validación externa con usuarios ni de una reevaluación formal con el instrumento heurístico del Capítulo 1 de este apartado. Ambas validaciones son, precisamente, el objeto de la fase siguiente.

\section{Fase 4 — Plan de validación (pendiente de ejecución)}

A la fecha de redacción de esta memoria, el prototipo \textbf{no ha sido validado} con usuarios reales ni reevaluado formalmente con el instrumento heurístico SISIB. Se cuenta, en cambio, con un plan de validación ya diseñado, que se ejecutará en el trabajo posterior a este capítulo y cuyos resultados se incorporarán en las secciones de Resultados, Discusión y Conclusiones. El plan contempla cinco fases:

\begin{enumerate}
    \item \textbf{Preparación del instrumento}: construir una rúbrica que replique las preguntas de chequeo del instrumento SISIB-UChile para las veinte dimensiones generales y los catorce factores específicos de selección escolar, con una hoja de cálculo que compute automáticamente el porcentaje de cumplimiento por dimensión y el puntaje global, siguiendo la misma fórmula del informe original (indicadores cumplidos sobre indicadores aplicables, por cien).
    \item \textbf{Aplicación al prototipo}: recorrer las secciones del prototipo aplicando cada indicador mediante cuatro tipos de evidencia —validadores automáticos (W3C, axe-core, Lighthouse, contraste WCAG), inspección visual en tres \emph{viewports} (375, 768 y 1280~px), inspección de código para indicadores estructurales, y pruebas funcionales manuales del simulador y del flujo de postulación completo—, incluyendo navegación exclusivamente por teclado y con lector de pantalla.
    \item \textbf{Análisis comparativo}: construir una tabla y gráficos que comparen el puntaje del prototipo contra el \emph{baseline} de la Fase~1, dimensión por dimensión, e identificar si las brechas críticas declaradas (inclusión, búsqueda, audiovisualidad, interoperabilidad) quedan efectivamente resueltas.
    \item \textbf{Redacción del informe}: documentar los hallazgos en un informe ejecutivo, con evidencia por indicador y un anexo de capturas de pantalla.
    \item \textbf{Triangulación con un segundo evaluador}: aplicar el mismo instrumento de forma independiente con una segunda persona evaluadora, y calcular un coeficiente de acuerdo entre evaluadores para mitigar el sesgo de autoevaluación.
\end{enumerate}

De manera complementaria, y siguiendo el precedente metodológico del estudio previo mencionado en la Sección~3.2, se planea una prueba de usabilidad con usuarios reales sobre el prototipo completo —no solo sobre el módulo de explicación del resultado—, replicando el esquema de tareas dirigidas y cuestionario de percepción en escala Likert utilizado en ese estudio previo, de modo que sus resultados sean comparables con los ya obtenidos para esa iteración anterior, además de evaluables contra el \emph{baseline} heurístico del sistema real.

Los resultados de esta fase, una vez ejecutada, permitirán responder con evidencia propia si el prototipo efectivamente mejora la comprensión, la confianza y la percepción de justicia de las familias respecto del proceso de admisión escolar —que es la pregunta de investigación planteada en el Capítulo~1— en lugar de inferirlo únicamente a partir del cumplimiento de una matriz de trazabilidad interna o de un estudio previo sobre un artefacto distinto.

\chapter{Resultados}

Este capítulo reporta únicamente resultados verificables de forma independiente a la fecha de esta memoria: el diagnóstico del SAE real, el estado objetivo de implementación del prototipo, y una comparación del prototipo con el estado actual del sitio oficial. \textbf{No se incluyen aquí resultados de validación con usuarios del prototipo desarrollado en este trabajo}, porque esa validación —descrita como plan en la Sección~3.5— aún no se ha ejecutado. Los resultados de un estudio de usabilidad previo, realizado sobre un artefacto distinto (una maqueta estática centrada solo en el módulo de explicación del resultado), se mencionan exclusivamente como antecedente en el Capítulo~3 y no se presentan en este capítulo como hallazgos propios.

\section{Diagnóstico heurístico del SAE real}

La evaluación heurística de calidad web aplicada al SAE real \citep{moralesvargas2026informe} arrojó un cumplimiento de \textbf{51\,\%} en el sitio informativo y de \textbf{61\,\%} en la plataforma de postulación, ambos con un desglose de 55\,\% en indicadores imprescindibles, 57\,\% en esperables y 22\,\% en deseables. Entre las brechas más críticas identificadas destacan:

\begin{itemize}
    \item \textbf{Transparencia y apertura} (0\,\% en el sitio informativo): no existe ninguna sección que explique públicamente el algoritmo de asignación, sus criterios o sus resultados agregados.
    \item \textbf{Inclusión} (0\,\%): sin opciones para necesidades educativas especiales ni lenguaje simplificado en el sitio informativo.
    \item \textbf{Búsqueda y encontrabilidad}: sin buscador interno ni mapa del sitio.
    \item \textbf{Audiovisualidad}: ausencia total de videotutoriales o infografías explicativas del proceso.
    \item \textbf{Interoperabilidad}: la contraseña de apoderado se crea de forma separada, sin integración con ClaveÚnica, pese a que esta ya contiene los datos requeridos en el registro (RUT, fecha de nacimiento, correo, teléfono).
    \item \textbf{Legibilidad}: el sitio informativo obtuvo un índice de legibilidad Spaulding de 120{,}43 (``difícil''), muy por sobre el nivel 6.\textsuperscript{o} básico recomendado para el público objetivo.
\end{itemize}

Este diagnóstico —detallado en veinte dimensiones para cada componente del SAE— constituye el \emph{baseline} cuantitativo que orientó las decisiones de diseño descritas en el Capítulo~3, y contra el cual debe compararse la reevaluación pendiente del prototipo (Sección~3.5).

\section{Estado de implementación del prototipo}

A partir del historial de desarrollo reconstruido en la Tabla~3.1, el prototipo evolucionó de un archivo HTML autocontenido de una sola página (Iteración~1) a una segunda versión HTML más completa (Iteración~2), y de ahí a una aplicación web construida con React, Vite y Tailwind CSS que cubre las cinco secciones exigidas por el diagnóstico original: página de inicio con buscador, módulo ``¿Cómo funciona el algoritmo?'' con simulador interactivo, ficha de colegio ampliada, flujo de postulación de tres pasos con inicio de sesión simulado vía ClaveÚnica, y panel de seguimiento de postulación con explicación contextualizada del resultado. Durante el desarrollo se agregaron además un comparador de colegios, un tour guiado de \emph{onboarding} y una página de ``Notas de diseño'' con la trazabilidad completa entre literatura y componentes, ninguno de ellos exigido explícitamente por el diagnóstico pero identificados como oportunidades de mejora durante el proceso iterativo.

Una verificación interna de trazabilidad frente a la matriz de problemas del diagnóstico (organizada en 20 categorías y aproximadamente 62 puntos de verificación) registra un cumplimiento de \textbf{62 de 62 puntos aplicables (100\,\%)} a nivel de código, incluyendo la totalidad de los puntos de prioridad alta y media; el último punto, una mejora de prioridad baja (sección de estadísticas históricas de asignación con datos ficticios), fue cerrado en julio de 2026. Un punto adicional de cabeceras HTTP de seguridad no aplica al prototipo por depender de la configuración de un servidor de producción.

Se insiste en que esta cifra es una \textbf{autoevaluación de implementación de código}, verificable revisando el propio prototipo, pero \textbf{no constituye} evidencia de que el prototipo mejore efectivamente la comprensión, la confianza o la percepción de justicia de usuarios reales: esa es precisamente la pregunta que el plan de validación de la Sección~3.5 busca responder, y que a la fecha de esta memoria permanece abierta.

\section{Estado actual del sitio oficial del SAE}

Como parte de esta investigación se realizó, adicionalmente, una revisión del estado vigente del sitio oficial del SAE (\texttt{sistemadeadmisionescolar.cl} y la plataforma de postulación en \texttt{admision.mineduc.cl}), con el fin de determinar si alguna de las brechas del diagnóstico original había sido cerrada de forma independiente por el Ministerio de Educación desde la evaluación de marzo de 2026. Se encontró que:

\begin{itemize}
    \item La \textbf{búsqueda y encontrabilidad} mejoró de forma concreta: existe ahora una ``Vitrina de Establecimientos'' con filtros por región, comuna, nivel educativo y nombre de establecimiento, con vistas de lista y mapa.
    \item La \textbf{inclusión} avanzó parcialmente: la vitrina indica explícitamente si un colegio cuenta con Programa de Integración Escolar (PIE), mediante un ícono dedicado.
    \item La \textbf{transparencia del algoritmo} mejoró solo de forma incremental: existen ahora páginas y preguntas frecuentes que describen los cuatro criterios de prioridad y desmienten mitos comunes (por ejemplo, que postular el primer día otorga ventaja), pero la explicación sigue siendo genérica y no contextualizada al caso particular del apoderado, y no existe ningún simulador interactivo.
    \item La \textbf{interoperabilidad con ClaveÚnica} no se resolvió: el acceso a la plataforma sigue requiriendo una cuenta separada con RUT/IPA y contraseña propia.
    \item La \textbf{audiovisualidad} tampoco se resolvió: no se encontraron videos ni infografías embebidas en el sitio informativo.
\end{itemize}

En síntesis, de las seis brechas prioritarias del diagnóstico original, el sitio oficial cerró de forma parcial dos (búsqueda e inclusión) y no ha resuelto las restantes cuatro al momento de esta revisión. Esta comparación es un resultado propio y verificable de esta memoria —a diferencia de la validación de usabilidad del prototipo, que permanece pendiente— y se retoma en el Capítulo~5 únicamente en la medida en que no depende de datos de usuarios aún no recolectados.

\section{Resultados pendientes}

Quedan explícitamente pendientes, y no se reportan en este capítulo: (a) la reevaluación del prototipo con el instrumento heurístico SISIB descrita en la Sección~3.5; (b) la prueba de usabilidad del prototipo completo con usuarios reales; y (c) cualquier comparación cuantitativa entre el desempeño del prototipo y el \emph{baseline} del SAE real que dependa de datos de usuarios. Estos resultados se incorporarán en una versión posterior de esta memoria, junto con la discusión y las conclusiones correspondientes (Capítulos~5 y~6).

\chapter{Discusión}

\emph{Capítulo pendiente de desarrollo.} La discusión completa de este trabajo depende de los resultados de validación descritos como pendientes en la Sección~3.5 y en la Sección~4.4: la reevaluación heurística del prototipo y la prueba de usabilidad con usuarios reales sobre la aplicación completa. Sin esos datos, no corresponde todavía discutir si el prototipo mejora efectivamente la comprensión, la confianza o la percepción de justicia de las familias, que es la pregunta de investigación central de esta memoria (Capítulo~1).

Lo que sí puede discutirse en este punto, por estar basado en resultados ya obtenidos y verificables (Capítulo~4), es lo siguiente.

\section{Por qué el sitio oficial cerró solo dos de las seis brechas prioritarias}

La comparación entre el prototipo y el estado actual del sitio oficial del SAE (Sección~4.3) muestra un patrón que vale la pena anotar como hipótesis a discutir con más evidencia una vez completada la validación: las brechas que el sitio oficial cerró de forma independiente —búsqueda y encontrabilidad, e inclusión parcial mediante el indicador de Programa de Integración Escolar— corresponden a mejoras principalmente informativas y de bajo acoplamiento interactivo, mientras que las que permanecen abiertas —transparencia algorítmica contextualizada, controles interactivos e interoperabilidad con ClaveÚnica— exigen precisamente el tipo de intervención que la literatura revisada en el Capítulo~2 identifica como más efectiva pero también más costosa de implementar: explicaciones locales personalizadas por caso, controles interactivos con retroalimentación inmediata \citep{kim2021recommender,feddersen2024forecasting}, e integración de sistemas de autenticación externos. Esta observación es preliminar y debe tratarse con cautela: compara un diagnóstico de marzo de 2026 con una revisión no sistemática del sitio realizada en julio de 2026, y no controla por otros cambios que el Ministerio de Educación pueda haber implementado fuera de las páginas revisadas.

\section{Qué preguntas deberá responder la discusión completa}

Una vez ejecutado el plan de validación de la Sección~3.5, este capítulo deberá abordar, como mínimo:

\begin{enumerate}
    \item Si el puntaje heurístico del prototipo, reevaluado con el mismo instrumento SISIB del diagnóstico original, efectivamente cierra las brechas críticas identificadas (transparencia, inclusión, búsqueda, audiovisualidad, interoperabilidad), y en qué magnitud respecto del \emph{baseline} de 51\,\%/61\,\%.
    \item Si los resultados de la prueba de usabilidad sobre el prototipo completo son consistentes con los del estudio previo mencionado en la Sección~3.2 —realizado sobre un artefacto más acotado—, en particular respecto de si la explicabilidad contextualizada mejora la comprensión y la confianza sin necesariamente resolver la percepción de justicia del resultado, tal como sugiere la literatura sobre gestión de expectativas \citep{springerwhittaker2019,springerwhittaker2020}.
    \item Qué brechas, si las hubiera, persisten pese al rediseño, y si esas brechas son atribuibles a limitaciones de diseño, a limitaciones inherentes a un prototipo de datos ficticios, o a aspectos del problema que exceden el alcance de una intervención de interfaz (por ejemplo, la interoperabilidad real con ClaveÚnica, que requiere integración de backend no simulable en un prototipo estático).
    \item Qué relación existe entre los hallazgos obtenidos y las condiciones de un despliegue real en producción, dado que la literatura revisada en el Capítulo~2 advierte explícitamente sobre la escasez de estudios de campo frente a estudios de laboratorio.
\end{enumerate}

\section{Limitaciones ya identificables}

Aun sin los resultados de validación, es posible anticipar limitaciones metodológicas que la discusión completa deberá considerar:

\begin{itemize}
    \item \textbf{Datos ficticios}: el catálogo de colegios y los perfiles de usuario del prototipo son ficticios pero representativos del contexto chileno; no se han usado datos reales de postulantes ni de establecimientos.
    \item \textbf{Persona única}: el diseño se validó conceptualmente contra un solo arquetipo de usuario (Daniela González); otros perfiles de apoderados —por ejemplo, familias de estudiantes con necesidades educativas especiales o usuarios migrantes— podrían revelar necesidades de diseño no cubiertas por este trabajo.
    \item \textbf{Alcance del algoritmo}: el prototipo no modifica ni audita el algoritmo de asignación real del SAE, cuyo código fuente no es público; el trabajo se limita a la capa de comunicación de sus resultados.
    \item \textbf{Comparación con el sitio oficial}: la revisión de la Sección~4.3 es una inspección puntual, no un nuevo diagnóstico heurístico completo con el instrumento SISIB; su alcance es más limitado que el del \emph{baseline} original.
\end{itemize}

\chapter{Conclusiones}\label{cap:conclusiones}

\emph{Capítulo pendiente de desarrollo.} Las conclusiones definitivas de esta memoria dependen de los resultados de validación descritos como pendientes en la Sección~\ref{sec:validacion} (reevaluación heurística del prototipo y prueba de usabilidad con usuarios reales), que a su vez determinan el contenido del Capítulo~\ref{cap:discusion}. Redactar conclusiones sobre la efectividad del prototipo antes de contar con esa evidencia sería anticipar un resultado no verificado, lo que esta memoria busca evitar de forma consistente con el propio principio de transparencia que defiende para el SAE.

\section*{Lo que puede afirmarse hasta ahora}
\addcontentsline{toc}{section}{Lo que puede afirmarse hasta ahora}

Sin adelantar conclusiones sobre efectividad, el trabajo realizado hasta la fecha permite sostener lo siguiente, con el respaldo de los Capítulos~\ref{cap:introduccion} a~\ref{cap:resultados}:

\begin{enumerate}
    \item La brecha de transparencia del SAE real está documentada con evidencia propia (diagnóstico heurístico, Sección~\ref{sec:res-diagnostico}) y no es una percepción subjetiva: el sitio informativo obtuvo 0\,\% de cumplimiento en transparencia y apertura, y 0\,\% en inclusión.
    \item Existe un cuerpo consistente de literatura científica (Capítulo~\ref{cap:marco}) que identifica principios de diseño concretos y transferibles —divulgación progresiva, explicabilidad contextualizada, controles interactivos con retroalimentación— para abordar brechas de este tipo en sistemas algorítmicos de diversos dominios.
    \item Es técnicamente posible construir un prototipo funcional que implemente esos principios sin modificar el algoritmo de asignación subyacente: el prototipo desarrollado en este trabajo aborda, a nivel de código, la totalidad de los puntos aplicables de la matriz de verificación derivada del diagnóstico original y sus ampliaciones posteriores (Sección~\ref{sec:res-prototipo}). Se trata de una autoevaluación de implementación —cuyo conteo agregado quedó pendiente de re-tabular con criterio explícito, según la nota de integridad de la Sección~\ref{sec:res-prototipo}—, no de una medición de cumplimiento validada externamente.
    \item El sitio oficial del SAE, dieciséis meses después del diagnóstico original, ha cerrado de forma independiente solo dos de las seis brechas prioritarias identificadas (Sección~\ref{sec:res-sitio-oficial}), lo que sugiere que el problema de investigación de esta memoria sigue vigente y no ha sido resuelto por otra vía.
\end{enumerate}

\section*{Lo que falta para concluir el trabajo}
\addcontentsline{toc}{section}{Lo que falta para concluir el trabajo}

En orden de prioridad:

\begin{enumerate}
    \item \textbf{Ejecutar la reevaluación heurística del prototipo} con el mismo instrumento SISIB usado en el diagnóstico original, siguiendo el plan de cinco fases descrito en la Sección~\ref{sec:validacion}, e idealmente con triangulación de un segundo evaluador.
    \item \textbf{Ejecutar la prueba de usabilidad con usuarios reales} sobre el prototipo completo (no solo sobre el módulo de explicación del resultado, a diferencia del estudio previo mencionado en la Sección~\ref{sec:fase2}), replicando un esquema de tareas dirigidas y cuestionario de percepción comparable.
    \item \textbf{Redactar el Capítulo~\ref{cap:discusion} (Discusión)} interpretando esos resultados a la luz de la literatura del Capítulo~\ref{cap:marco} y del estado del sitio oficial reportado en la Sección~\ref{sec:res-sitio-oficial}.
    \item \textbf{Redactar este capítulo de Conclusiones} en su versión definitiva, incluyendo las contribuciones reales del trabajo, sus limitaciones confirmadas empíricamente, y las líneas de trabajo futuro —entre ellas, la colaboración con el Ministerio de Educación y la ampliación a otros perfiles de usuario— que ya se perfilan en la Sección~\ref{sec:limitaciones} pero que solo pueden fijarse con propiedad una vez conocidos los resultados.
\end{enumerate}


% ============================================================
%  BIBLIOGRAFÍA
% ============================================================
\bibliography{referencias}

\end{document}
